\section{A favorable physical filter and the initial vector}

\begin{lemma}[A favorable filter for the clipped ground vector]
    \label{thm:aklv_favorable_physical_filter}
    \lean{TNLean.MPS.AreaLaw.exists_aklv_physical_favorable_filter}
    \leanok
    \uses{thm:aklv_favorable_scalar_parameters,cor:aklv_logarithmic_clipped_filter}
    There is a universal constant $A\ge21$ with the following property.
    Let $H=\sum_{i=0}^{L-2}h_i$, where $0\le h_i\le\Id$ acts on
    $\{i,i+1\}$. Suppose $\Omega$ is a unit eigenvector of energy $\mu$,
    $g>0$, $\mu+g\in\operatorname{spec}(H)$, and
    \begin{align*}
        H-\mu\Id\ge g(\Id-|\Omega\rangle\langle\Omega|).
    \end{align*}
    Set
    \begin{align*}
        u&=1+\log q+\log(1/g),&
        s&=\lceil Au^2/g\rceil,&\ell&=s^2,\\
        t_0&=1024+16\log(1/g),&R&=2t_0+s,\\
        D&=\bigl(q(s^2+1)\bigr)^{32(2s+1)},&
        \Delta&=4\exp\left(-4\ell\sqrt{\frac{9g}{10R}}\right).
    \end{align*}
    Suppose the window fits: $a+s+1\le L$.
    Then $q\ge2$, $g\le1$, $s\ge21$, $\Delta>0$,
    $D\Delta\le10^{-5}$, and
    \begin{align*}
        \log D\le198(A+1)\bigl(6+2\log(A+1)\bigr)u^3/g.
    \end{align*}
    Choose the three actual reference minima and the centered ground energy
    $E_0\in[0,6]$ as in Lemma~\ref{cor:aklv_logarithmic_clipped_filter}.
    The actual two-exterior truncation $B$ at $t_0$ has a unit ground vector
    $\phi$, with energy $E\ge0$, such that
    \begin{align*}
        0\le B&\le R\Id,&
        B-E\Id&\ge\frac{9g}{10}(\Id-|\phi\rangle\langle\phi|),\\
        \|\phi-\Omega\|^2&\le e^{-t_0/16}/100.
    \end{align*}
    There is a real polynomial $p$ of degree $\ell$, with $p(E)=1$,
    such that $K=p(B)$ fixes $\phi$, preserves $\phi^\perp$, and satisfies
    \begin{align*}
        \operatorname{OSR}_{a+\lfloor s/2\rfloor}(K)&\le D,&
        \|K-|\phi\rangle\langle\phi|\|&\le\sqrt\Delta.
    \end{align*}
    Every unit $x\perp\phi$ satisfies $\|Kx\|^2\le\Delta$.
    This is the favorable initial filter used in the construction following
    \cite[Theorem~6.1]{Arad2013AreaLaw}, for the stated available-window range.
\end{lemma}
\begin{proof}
    \leanok
    The two distinct spectral values give $q\ge2$. The endpoint trial-vector
    argument for the original chain gives $g\le1$.
    Apply Lemma~\ref{thm:aklv_favorable_scalar_parameters} and then
    Lemma~\ref{cor:aklv_logarithmic_clipped_filter} at $t_0$ and
    $\ell=s^2$. The identities $\lfloor\ell/s\rfloor=s$ and
    $\lfloor\sqrt\ell\rfloor=s$ give the displayed rank bound.
    Squaring the amplitude estimate gives the stated $\Delta$.
\end{proof}

\begin{lemma}[The initial filtered product vector]
    \label{thm:aklv_physical_product_initialization}
    \lean{TNLean.MPS.AreaLaw.exists_aklv_physical_product_initialization}
    \leanok
    \uses{thm:aklv_favorable_physical_filter}
    Under the hypotheses and with the parameters of
    Lemma~\ref{thm:aklv_favorable_physical_filter}, the regional minima,
    clipped ground vector $\phi$, and polynomial filter $K$ can be chosen
    together with a unit product vector $\chi=v_L\otimes v_R$ across
    the cut $a+\lfloor s/2\rfloor$ and a scalar $c$, $|c|=1$, such that
    \begin{align*}
        |\langle\phi,\chi\rangle|&\ge(2D)^{-1/2},&K\chi&\ne0.
    \end{align*}
    The actual normalized output $\psi_0=cK\chi/\|K\chi\|$ satisfies
    \begin{align*}
        \|\psi_0\|&=1,&\operatorname{SR}(\psi_0)&\le D,\\
        \|\psi_0-\phi\|^2&\le1/64,&
        \|\psi_0-\Omega\|^2&\le1/16.
    \end{align*}
    The overlap estimate concerns the clipped target $\phi$.
    This proves the initial step of \cite[Lemma~6.3]{Arad2013AreaLaw}
    for the stated available-window range.
\end{lemma}
\begin{proof}
    \leanok
    Apply the product-overlap argument of
    \cite[Lemma~2.2]{Arad2013AreaLaw} to the fixed vector $\phi$ of $K$;
    the required inequality $D\Delta\le1/2$ follows from
    $D\Delta\le10^{-5}$. Write $b=\langle\phi,\chi\rangle$ and
    $z=\chi-b\phi$. Orthogonal projection gives $\|z\|\le1$.
    Choosing the phase of the normalized output gives
    \begin{align*}
        \|\psi_0-\phi\|^2
        &\le\frac{2\Delta\|z\|^2}{|b|^2}
          \le4D\Delta\le1/64.
    \end{align*}
    Multiplication by $K$ increases the Schmidt rank of the product vector
    by at most its operator Schmidt rank, and multiplication by a nonzero
    scalar preserves that rank. Finally $t_0\ge0$, so
    \begin{align*}
        \|\phi-\Omega\|^2&\le1/100,\\
        \|\psi_0-\Omega\|^2
        &\le2\|\psi_0-\phi\|^2+2\|\phi-\Omega\|^2\le1/16.
    \end{align*}
\end{proof}
